% =============================================================================
% Project : THE_PHILOSOPHY_OF_VIABILITY
% Book    : Book I — Ontology of Viability
% Book ID : VS.B01
%
% File    : 4_ASSETS/tables/VS.B01.P01.C03.S08.SS13.P003.TL01.methodological_limits_of_unity_and_plurality.tex
%
% Asset:
%   Table TL01 — Methodological Limits of Unity and Plurality
%
% Purpose :
%   To summarize the principal restrictions governing claims
%   about unity, plurality, and their implications for identity.
%
% =============================================================================

\begin{table}[htbp]
\centering
\caption{Methodological Limits of Unity and Plurality}
\label{tab:methodological-limits-of-unity-and-plurality}
\footnotesize
\renewcommand{\arraystretch}{1.2}

\begin{tabularx}{\textwidth}{
    @{}
    p{0.24\textwidth}
    X
    p{0.32\textwidth}
    @{}
}
\toprule
\textbf{Analytical Requirement} &
\textbf{Restriction} &
\textbf{Inference Not Licensed by Itself}
\\
\midrule

Specify the subject &
Identify what is unified, plural, counted, or distinguished. &
Unity of one subject establishes unity of another subject.
\\

Specify the respect &
Identify the ontological determination under examination. &
Material, relational, or manifestational plurality establishes identity plurality.
\\

Specify the domain &
Identify the range within which the claim is valid. &
A local numerical result applies to every ontological domain.
\\

Specify the criterion &
Identify what counts as one or multiple. &
Numerical classification determines ontological identity-number.
\\

Preserve conceptual distinctions &
Keep source, aggregate, composition, structure, manifestation, and identity distinct. &
Unity or plurality in one category settles another category.
\\

Establish cross-respect relations &
Provide grounds for transferring a conclusion between respects. &
Co-occurrence of determinations establishes entailment.
\\

Preserve unresolved questions &
Do not convert insufficient evidence into a universal conclusion. &
Failure to prove plurality establishes singular identity, or vice versa.
\\

\bottomrule
\end{tabularx}
\end{table}

% =============================================================================
% END OF TABLE TL01
% =============================================================================